\subsection{Characteristic of a ring and of a field}

We shall define the characteristic of a ring A only when A has a subfield. When this is so, let f be the unique ring homomorphism of Z into A, and let n be the unique positive integer generating the ideal of Z which is the kernel of f (I, p.~111); then the integer n is called the characteristic of A.

Let A be a ring for which the characteristic is defined ; then A does not reduce to 0. By Th. 1 there exists a unique subfield P of A which is a prime field ; we shall call it the prime subfield of A. By the proof of Th. 1 there are the following two possibilities :

\begin{enumerate}
\def\labelenumi{\alph{enumi})}
\item
  the characteristic of A is 0, P is isomorphic to Q,
\item
  the characteristic of \(\mathbf{A}\) is a prime number \(p\) , \(\mathbf{P}\) is isomorphic to \(\mathbf{F}_p\) .
\end{enumerate}

If the characteristic of A is zero, there exists a unique ring homomorphism of Q into A; its image is the prime subfield of A, contained in the centre of A. Therefore there exists a unique Q-algebra structure of A compatible with the ring structure. When the characteristic of A is a prime number p, we have the corresponding properties on replacing the field Q by the field \(F_{p}\) .

PROPOSITION 1. --- Let A be a ring not reduced to 0.

\begin{enumerate}
\def\labelenumi{\alph{enumi})}
\item
  For A to be of characteristic 0 it is necessary and sufficient that the mapping \(x \mapsto n \cdot x\) of A into itself should be bijective, for every integer \(n \neq 0\) .
\item
  Let \(p\) be a prime number. For \(\mathbf{A}\) to be of characteristic \(p\) it is necessary and sufficient that \(p \cdot x = 0\) for all \(x \in \mathbf{A}\) .
\end{enumerate}

Let \(f\) be the unique homomorphism of \(\mathbf{Z}\) into \(\mathbf{A}\) ; we have \(n \cdot x = f(n)x\) for every integer \(n\) and every \(x\) in \(\mathbf{A}\) . For \(\mathbf{A}\) to be of characteristic 0, it is necessary and sufficient that \(f\) should extend to a homomorphism of \(\mathbf{Q}\) into \(\mathbf{A}\) , that is, \(f(n)\) should be invertible in \(\mathbf{A}\) for every \(n \neq 0\) (I, p.~113); this proves \(a\) ). Similarly for \(\mathbf{A}\) to be of characteristic \(p\) it is necessary and sufficient that \(f\) should annihilate \(p\mathbf{Z}\) , that is, \(f(p) = 0\) , or also that \(p \cdot x = 0\) for all \(x \in \mathbf{A}\) ; this proves \(b\) .

Let us take for A a not necessarily commutative field. The centre of A is a (commutative) field; therefore the characteristic and the prime subfield of A are defined.

Remarks. --- 1) Let A and A' be two rings not reduced to 0. Suppose that the characteristic of A is defined and that there is a homomorphism u of A into A'. The image under u of the prime subfield of A is a subfield P' of A', isomorphic to P, and hence prime. It follows that the characteristic of A' is defined and is equal to that of A. If A and A' are of characteristic 0 (resp. \(p \neq 0\) ), the mapping u is a homomorphism of algebras over Q (resp. \(F_{p}\) ).

\begin{enumerate}
\def\labelenumi{\arabic{enumi})}
\setcounter{enumi}{1}
\item
  Remark 1 shows that if A is a ring of characteristic 0 (resp. \(p \neq 0\) ), the same holds of any ring A' containing A as subring, or of any quotient of A by a two-sided ideal \(\mathfrak{a} \neq \mathbf{A}\) . In particular, if K is a field, every subfield of K and every extension field of K have the same characteristic as K.
\item
  Let A be an algebra not reduced to 0 over a field K. Since the mapping \(\lambda \mapsto \lambda \cdot 1\) of K into A is a ring homomorphism, Remark 1 shows that the characteristic of A is defined and equal to that of K.
\item
  Since the field \(\mathbf{Q}\) is infinite, every ring of characteristic 0 is infinite; it follows that every finite field has non-zero characteristic.
\item
  Let A be a ring not reduced to 0, whose additive group is a torsion-free Z-module, and put \(B = Q \otimes_{Z} A\) . The mapping \(x \mapsto 1 \otimes x\) of A into B is injective (II, p.~314), hence A is isomorphic to a subring of a ring of characteristic 0.
\end{enumerate}

